\section{Related Work}

\textbf{Accessibility Studies.} \citet{9555469} found that BLV and sighted reader groups differ significantly on which semantic content they consider as most useful, suggesting that access to meaningful information is strongly reader-specific. VizWiz-VQA \cite{gurari2018vizwiz} contains images and visual QA pairs produced by blind people encouraging the development of more generalized algorithms that can assist the blind. As an extended work, VizWiz-LF \cite{huh2024long} includes long-form answers from BLV people. VisText \cite{tang2023vistext} contains charts and captions that convey different levels of semantic content. As shown in Table~\ref{tab:dataset_comparison}, VizWiz-VQA and VizWiz-LF were validated by BLV users but only focus on Visual QA (VQA) applications. VisText examines the role of the level of semantic content but was not validated by BLV for dataset purposes. As a diagram description dataset validated by BLV users, \textsc{Sightation} explores diverse use cases, with assessments on various aspects. 

\noindent \textbf{Image Description Tasks and Models.} \citet{wang2024qwen2} presented the \textsc{Qwen2-VL} collection, which includes three open-weights models: 2B, 7B, and 72B. \textsc{Qwen2-VL} matches the performance of \textsc{GPT-4o} and \textsc{Claude3.5-Sonnet} \cite{anthropic2024claude} in multimodal scenarios, surpassing other open-weights VLMs at the time. 

\textsc{GPT-4o} \cite{hurst2024gpt} accepts multimodal input and generates high-quality outputs including text and codes, showing powerful multimodal understanding capability. Using these VLMs, the image description task aims to generate a descriptive textual context for images of different types (\textit{e.g.}, photographs, illustrations, schematics, and diagrams). Flickr8K and PASCAL-50S comprise natural images, captions, and human judgments\cite{hodosh2013framing, vedantam2015cider}, and Polaris \cite{wada2024polos} incorporated synthetic captions from image captioning models.

ChartGemma \cite{masry2024chartgemma} contains chart images collected from specialized websites and instruction-tuning data generated from the charts. MathVista \cite{lu2023mathvista} encompasses diverse visual contexts from natural images to diagrams or plots that require mathematical reasoning. However, Table~\ref{tab:dataset_comparison} shows that these datasets have an average text length much shorter than ours, even though charts and mathematical images could be highly information-dense. Complementing the limitation, \textsc{Sightation} provides contexts that top in average text length to date with variants for downstream tasks.

\noindent \textbf{Human Annotation Efforts.} 
Human judgment annotations are essential in evaluating image captions, complementary to automatic metrics. Common approaches involve employing annotators to assess captions based on rating scales for specific dimensions of text quality \cite{gehrmann2023repairing}. However, it comes with challenges, including subjectivity and consistency issues. \citet{amidei-etal-2019-agreement} argues that the evaluation of generated text is intrinsically subjective and relies on different factors including annotator experience, motivation, knowledge, or education.  A related line of research \citep{glockner-etal-2024-ambifc, nie2020can} directly addressing this limitation advocates that generations from few-annotator pools fall short in terms of coverage of the distribution of opinions.
